\documentclass[12pt]{article}
\usepackage[utf8]{inputenc}
\usepackage[T1]{fontenc}
\usepackage{hyperref}
\title{Tatar Language Resources for K-12 Curriculum Development}
\author{Compiled by Claude Sonnet 5 (claude-sonnet-5) for Andr\'as Kornai}
\date{}
\begin{document}
\maketitle

\section{Introduction}
Tatar (ISO 639-3 \texttt{tat}, Glottocode tata1255) is a Turkic language with about 5.4 million speakers and is an official language in Tatarstan, Russia. This report is based on a single web search. It surfaced policy news, academic papers on lexicography and textbooks, and corpus descriptions, but not the textbooks, dictionaries or corpus sites themselves. Much more exists than was found here.

\section{Primary Educational Resources}
\begin{itemize}
\item Tatar World Congress news item: the Ministry of Education and Science of Tatarstan will approve printing of new Tatar-language textbooks. The page reports 89 Tatar language and literature textbooks on the federal list, new programmes for Tatar as state language, as mother tongue and for Tatar literature, and a pilot of the new federal standard in 50 schools. The date is not shown. \url{https://tatar-congress.org/en/news/printing-of-new-textbooks-on-the-tatar-language-will-be-approved-by-the-ministry-of-education-and-science-of-tatarstan/}
\item Related Tatar World Congress page from the same search. \url{https://tatar-congress.org/en/?p=323542}
\item Kazan Federal University paper on teaching Tatar in Russian schools, recommending textbooks for learners of Tatar as a foreign language. \url{https://dspace.kpfu.ru/xmlui/handle/net/138774}
\item Kazan Federal University study analysing the structure and content of Tatar textbooks for non-Russian students. \url{https://dspace.kpfu.ru/xmlui/handle/net/140428}
\end{itemize}

\section{Linguistic and Reference Materials}
\begin{itemize}
\item Gulshat Safiullina, paper on the history of Tatar lexicography, from terminological works to learner and picture dictionaries, noting the lack of English--Tatar dictionaries (EURALEX 2016). \url{https://euralex.org/wp-content/themes/euralex/proceedings/Euralex%202016/euralex_2016_052_p475.pdf}
\item EURALEX record for ``Bilingual lexicography in the Republic of Tatarstan in 1990--2010''. \url{https://euralex.org/publications/bilingual-lexicography-in-the-republic-of-tatarstan-in-1990-2010/}
\item EURALEX record for the Safiullina paper. \url{https://euralex.org/?p=7759}
\end{itemize}

\section{Religious and Community Resources}
The World Congress of Tatars (tatar-congress.org) appears as a community organisation source above. No religious materials were found.

\section{Digital and Audio Resources}
\begin{itemize}
\item Kazan Federal University repository record describing the Tatar corpora (Corpus of Written Tatar at KFU and the Tatar National Corpus ``Tugan Tel''), as reported in the search summary. \url{https://dspace.kpfu.ru/xmlui/handle/net/142385}
\item Conference paper record on Tatar corpus-based grammar research. \url{https://library.iated.org/view/GALIEVA2017DEV}
\end{itemize}
I did not see the corpus sites' own URLs in the results.

\section{Curriculum Development Context}
Tatar is taught as a subject and as a state or mother tongue in Tatarstan, with federal-listed textbooks and new programmes. Gaps noted in the literature include English--Tatar dictionaries and English textbooks with Tatar-language instructions. Specific textbook titles and publishers should be taken from the official federal list rather than from this report.

\end{document}